\documentclass[12pt]{article}
\usepackage[utf8]{inputenc}
\usepackage[T1]{fontenc}
\usepackage{hyperref}
\title{Wumeng Nasu Language Resources for K-12 Curriculum Development}
\author{Compiled by Claude Sonnet 5.5 (claude-sonnet-5-5) for Andr\'as Kornai}
\date{}
\begin{document}
\maketitle

\section{Introduction}
Wumeng Nasu (ISO 639-3 ywu, Glottocode wume1235) is a Sino-Tibetan language of the Yi group, spoken mainly in Guizhou with some speakers in Yunnan. Joshua Project gives about 320,000 speakers for Nasu, Wumeng. The code history is tangled: an older code, ywm, was retired. The mapping should be checked against the current ISO 639-3 registry. Very little was found, and nothing aimed at schools.

\section{Primary Educational Resources}
No primers, textbooks or school programmes were found in the search results.

\section{Linguistic and Reference Materials}
\begin{itemize}
\item Wikipedia, Nasu language, which lists ywu among the Nasu codes and notes the Pollard script for Nasu: \url{https://en.wikipedia.org/wiki/Nasu_language}
\item Glottolog record for the retired code ywm (Wumeng Yi): \url{https://glottolog.org/resource/languoid/iso/ywm}
\end{itemize}

\section{Religious and Community Resources}
\begin{itemize}
\item Joshua Project profile (a related Wusa Nasu group; the match to ywu is uncertain): \url{https://joshuaproject.net/people_groups/18629}
\item Asia Harvest profile of the Wumeng people in Yunnan: \url{https://people-groups.asiaharvest.org/people-group-profiles/yunnan/wumeng.pdf}
\end{itemize}

\section{Digital and Audio Resources}
No audio or digital language resources specific to Wumeng Nasu were found.

\section{Curriculum Development Context}
No evidence was found that the language is taught in schools. Wumeng is related to, but only partly intelligible with, Shuixi, Mangbu and Wusa Nasu, so materials for those varieties should not be assumed to fit it. A curriculum would need new documentation work.
\end{document}